% ============================================================
% CHƯƠNG 2: CƠ SỞ LÝ THUYẾT
% ============================================================
\chapter{Cơ sở lý thuyết}

\section{Mô hình ngôn ngữ lớn (Large Language Models)}

\subsection{Kiến trúc Transformer}
% TODO:
% - Self-Attention mechanism
% - Multi-head attention
% - Encoder-Decoder architecture
% - Positional encoding
% Hình 2.1: Kiến trúc Transformer

\subsection{Các mô hình tiền huấn luyện}
% TODO:
% - BERT và các biến thể (encoder-only)
% - GPT và các biến thể (decoder-only)
% - T5 và các biến thể (encoder-decoder)
% - So sánh kiến trúc và ứng dụng phù hợp

\subsection{Các mô hình sử dụng trong khóa luận}
% TODO:
% 2.1.3.1. GPT-4o (OpenAI)
%   - Kiến trúc decoder-only, đa ngôn ngữ
%   - Khả năng zero-shot summarization
%   - Structured output với JSON mode
%
% 2.1.3.2. ViT5 (VietAI/vit5-large-vietnews-summarization)
%   - Kiến trúc encoder-decoder (T5-based)
%   - Pre-trained trên CC100-vi, fine-tuned trên VietnewsQA
%   - Chuyên biệt cho tóm tắt tin tức tiếng Việt
%
% 2.1.3.3. PhoBERT (vinai/phobert-base)
%   - Kiến trúc encoder-only (BERT-based)
%   - Pre-trained trên 20GB tiếng Việt
%   - Sử dụng cho BERTScore (đánh giá ngữ nghĩa)
%
% 2.1.3.4. PhoGPT-4B-Chat (vinai/PhoGPT-4B-Chat)
%   - Kiến trúc decoder-only (GPT-based), 4B tham số
%   - Thảo luận về thách thức triển khai (GPU requirements)

\section{Các phương pháp đánh giá chất lượng tóm tắt}

\subsection{ROUGE (Recall-Oriented Understudy for Gisting Evaluation)}
% TODO:
% - ROUGE-1 (unigram), ROUGE-2 (bigram), ROUGE-L (longest common subsequence)
% - Công thức tính Precision, Recall, F1
% - Ưu điểm: đơn giản, phổ biến
% - Nhược điểm: chỉ so sánh bề mặt từ vựng

\subsection{BLEU (Bilingual Evaluation Understudy)}
% TODO:
% - Nguyên lý: precision-based n-gram matching
% - Brevity penalty
% - Ưu điểm: standard metric
% - Nhược điểm: không phản ánh ngữ nghĩa

\subsection{BERTScore}
% TODO:
% - Nguyên lý: cosine similarity giữa các token embeddings
% - Sử dụng PhoBERT cho tiếng Việt (thay vì BERT tiếng Anh)
% - Ưu điểm: đánh giá ngữ nghĩa, phù hợp tóm tắt abstractive
% - Truncation 256 tokens cho PhoBERT
% Hình 2.2: Minh họa BERTScore

\subsection{Tỷ lệ nén (Compression Rate)}
% TODO:
% - Công thức: len(summary) / len(original)
% - Ý nghĩa: đánh giá mức độ cô đọng thông tin

\section{Kỹ thuật trích xuất nội dung web}

\subsection{Mozilla Readability}
% TODO:
% - Thuật toán phân tích DOM
% - Loại bỏ noise (quảng cáo, menu, sidebar)
% - Trích xuất nội dung chính

\subsection{Xử lý đặc thù các trang báo Việt Nam}
% TODO:
% - Cấu trúc DOM khác nhau giữa VnExpress, Tuổi Trẻ, Dân Trí, Thanh Niên
% - Xử lý ảnh, video, bảng biểu trong bài viết

\section{Kiến trúc Search-Augmented Generation}

\subsection{Retrieval-Augmented Generation (RAG)}
% TODO:
% - Nguyên lý: kết hợp tìm kiếm với sinh văn bản
% - Ứng dụng trong fact-checking: tìm nguồn xác minh

\subsection{Tavily Search API}
% TODO:
% - API tìm kiếm tối ưu cho AI/LLM
% - Trả về nội dung có cấu trúc, phù hợp làm context cho LLM

\section{Công nghệ phát triển tiện ích trình duyệt}

\subsection{Manifest V3 và Plasmo Framework}
% TODO:
% - Chrome Extension Manifest V3
% - Plasmo: framework React cho browser extension
% - Shadow DOM và content script injection

\section{Tổng kết chương}
% TODO: Tóm tắt các cơ sở lý thuyết đã trình bày
